\subsection{Independent verification of the cote 137 test objects}
\label{sec:cosimplicial-recheck}

We check the three supplied statements from RESEARCH\_SA09 that are
offered as additional tools for cote 136.  Their mathematical content
is considered directly; the supplied transcriptions of cote 137 are
not additional hypotheses in the arguments below.  We write
$\mathcal P=\operatorname{Ch}^{b}_{\geq0}
(\operatorname{Free}_{\mathrm{fg}}R)$, with degreewise split exact
sequences, and put $D_i=(R\xrightarrow{1}R)$ in homological degrees
$i+1,i$.  Its map $u_i:D_i\to D_{i+1}$ is the identity from the upper
term of the first disk to the lower term of the second.

\paragraph{A-S09-01 (Proved; independent re-proof of SA09 A2.2).}
Let $R$ be any commutative ring, including the zero ring, and let
$\mathcal B$ be an $R$-linear additive category.  The disk restriction
of relatively left-exact additive $R$-linear functors
$\mathcal P\to\mathcal B$ is fully faithful.  Its objects are precisely
the nonnegative cochain complexes $C$ for which the following kernels
exist for every $L\in\mathcal P$:
\[
 T_C(L)=\ker\left[
 \bigoplus_i L_i\otimes_R C^i\longrightarrow
 \bigoplus_i L_i\otimes_R C^{i+1}\right],
 \qquad
 (x_i)\longmapsto(d x_i-\partial_{i+1}x_{i+1})_i .
 \tag{S09.1}
\]
All sums are finite; tensoring by a finite free $R$-module means its
finite copower in $\mathcal B$.

\begin{proof}
For $b\in\mathcal B$, the universal property of a kernel identifies
\[
 \operatorname{Hom}_{\mathcal B}(b,T_C(L))
 =\operatorname{Hom}_{\operatorname{Ch}(R)}
 \bigl(L^\vee,\operatorname{Hom}_{\mathcal B}(b,C)\bigr),
\]
where $(L^\vee)^i=\operatorname{Hom}_R(L_i,R)$ and its differential
is $\partial_{i+1}^{\vee}$.  Dualizing a degreewise split exact
sequence reverses it and remains exact.  Left exactness of the Hom
functor proves the required relative left exactness of $T_C$.
The same formula proves functoriality in $L$ and in $C$.  A map
$D_i^\vee\to C$ is determined by its degree-$i$ component, hence
$T_C(D_i)=C^i$; the map induced by $u_i$ is $d^i$.

For the converse set $I(L)=\bigoplus_i D_i\otimes_RL_i$.
Its degree-$n$ term is $L_n\oplus L_{n-1}$, with differential
$(x,y)\mapsto(y,0)$.  The maps
\[
 j_L(x)=(x,\partial x),\qquad q_L(x,y)=y-\partial x
\]
give a degreewise split exact sequence
$0\to L\xrightarrow{j_L}I(L)\xrightarrow{q_L}Q(L)\to0$,
where $Q(L)_n=L_{n-1}$ and $d_Q=-\partial$.
Applying the given functor $F$ to this sequence and to the analogous
sequence for $Q(L)$ identifies $F(L)$ with the kernel of
$F(j_{Q(L)}q_L):F(I(L))\to F(I(Q(L)))$.
Its diagonal disk-to-next-disk entries are $F(u_i)$ and its
same-disk entries are $-\partial_{i+1}$, exactly the matrix in
(S09.1).  This supplies a natural isomorphism $F\cong T_{\rho F}$.
A cochain map on disk values induces maps on all the kernels.
Conversely, naturality with $j_L$, $q_L$ and their disk copresentations
forces any natural transformation to be this induced map.  This proves
full faithfulness and the stated essential image.  No exactness of
tensoring with an arbitrary $R$-module has been used.
\end{proof}

\paragraph{A-S09-02 (Proved; independent re-proof of SA09 A3.1).}
Let $\mathcal T$ be the category of simplicial abelian groups with
bounded finite-free normalization.  For $n\geq0$ define
\[
 (R_n)_m=\mathbb Z^{\operatorname{Hom}_\Delta([n],[m])},
\]
with simplicial maps given by restriction of functions.
Then $R_n\in\mathcal T$ and has an integral normalized contraction.
In the normalization $N_m=\bigcap_{i=1}^{m}\ker d_i$, the rank is
$\binom{n+1}{m}$ for $0\leq m\leq n+1$ and is zero otherwise.
The cosimplicial object $R_\bullet$ has cosimplicial normalization
$(D_i,u_i)$.  The category $\mathcal T$ is closed under levelwise
tensor products and all divided powers $\Gamma^p$, including $p=0$.

\begin{proof}
Index the characteristic-function basis in degree $m$ by
nondecreasing tuples $t=(t_0,\ldots,t_n)$ with entries in $[m]$.
Restriction along $d_i$ selects exactly the tuples avoiding $i$.
Thus the normalized basis consists of the tuples containing every
integer $1,\ldots,m$.  The multiplicities of these $m$ integers are
positive, while that of zero is nonnegative.  Counting their sum
$n+1$ gives $\binom{n+1}{m}$.

Writing $e_t$ for a normalized basis vector, the normalized
differential and a contraction are
\[
 d(e_t)=
 \begin{cases}e_{t-1}&0\notin t,\\0&0\in t,\end{cases}
 \qquad
 h(e_t)=
 \begin{cases}e_{t+1}&0\in t,\\0&0\notin t.\end{cases}
 \tag{S09.2}
\]
Here $d=0$ out of degree zero.  When $0\notin t$, the tuple $t-1$
contains zero and all positive integers required in the target.
When $0\in t$, the tuple $t+1$ contains all required positive
integers in the next degree.  Exactly one of $dh(e_t)$ and $hd(e_t)$
equals $e_t$, and the other is zero.  Hence $dh+hd=1$ over
$\mathbb Z$, with $d^2=h^2=0$.  This proves split acyclicity,
including degree zero.  The attached integer certificate checks
every basis vector for $0\leq n\leq8$, including two degrees beyond
the support: 1,022 basis vectors and 2,044 individual identities.

For $L\in\mathcal T$, the right-adjoint evaluation formula gives
$\operatorname{Hom}_{s\mathsf{Ab}}(L,R_n)
=\operatorname{Hom}_{\mathbb Z}(L_n,\mathbb Z)$.
Cosimplicial normalization on the right gives
$\operatorname{Hom}_{\mathbb Z}((NL)_i,\mathbb Z)$ with differential
dual to $d_L$.  This functor is represented in $\mathcal T$ by the
simplicial denormalization of $D_i$.  Yoneda identifies the normalized
objects and their differential with $D_i$ and $u_i$.

Finally, for $L\in\mathcal T$ whose normalization vanishes above
degree $a$, the Dold--Kan splitting gives
$\operatorname{rank}L_m=\sum_{i=0}^{a}\binom mi
\operatorname{rank}(NL)_i$.
The product of two such rank polynomials has degree at most $a+b$;
the rank polynomial of $\Gamma^pL$ has degree at most $pa$ for $p>0$.
All levels are finite free, and normalization is a direct summand in
each level.  Uniqueness of expansion in the binomial-polynomial basis
forces the normalized terms to vanish above these respective bounds.
For $p=0$ the object is the constant group $\mathbb Z$.
\end{proof}

\paragraph{A-S09-03 (Proved; targeted re-proof of SA09 A3.2).}
For every simplicial set $X$, including infinite ones, set
$F_X(M)=\operatorname{Hom}_{s\mathsf{Ab}}(\mathbb Z[X],M)$ on
$\mathcal T$.  Its full natural action
\[
 \rho_{p,M}:\Gamma^p F_X(M)\longrightarrow F_X(\Gamma^pM)
\]
is defined by applying $\Gamma^p$ to evaluation at every simplex.
The integral normalized cochain complex together with this action
determines, and is determined by, the cosimplicial function ring
$A_X^n=\operatorname{Map}(X_n,\mathbb Z)$.  The reconstruction is
on the image of $X\mapsto F_X$, with natural isomorphisms of the
specified data.

\begin{proof}
The underlying additive functor is relatively left exact: applying
$\operatorname{Hom}(\mathbb Z[X],-)$ to a short exact sequence of
simplicial groups gives its kernel diagram.  Its disk values are
$\operatorname{Hom}_{\mathbb Z}((N\mathbb Z[X])_i,\mathbb Z)$.
A-S09-01 therefore reconstructs this functor from the normalized
cochains without a finiteness restriction on $X$.

The middle direct summand in
$\Gamma^2(P\oplus Q)=\Gamma^2P\oplus(P\otimes Q)\oplus\Gamma^2Q$
recovers from $\rho_2$ the map
$F_X(M)\otimes F_X(N)\to F_X(M\otimes N)$ sending $(f,g)$ to
the pointwise function $x\mapsto f(x)\otimes g(x)$.
The map $\rho_0$ is the constant-function unit.  Evaluation adjunction
identifies $F_X(R_n)$ with $\mathbb Z^{X_n}$; the pointwise ring
multiplication $R_n\otimes R_n\to R_n$, composed with the recovered
pairing, gives the multiplication on this function ring.  The
cosimplicial operators follow from the maps $R_n\to R_{n'}$ of
A-S09-02.

Conversely $M_m$ is finite free, so the map
$A_X^m\otimes M_m\to\operatorname{Map}(X_m,M_m)$ sending
$a\otimes v$ to $(x\mapsto a(x)v)$ is an isomorphism even for
infinite $X_m$.  The equalizer
\[
 \prod_{m\geq0} A_X^m\otimes M_m
 \ \rightrightarrows\
 \prod_{\theta:[m]\to[n]} A_X^n\otimes M_m
 \tag{S09.3}
\]
recovers $F_X(M)$: its two arrows on the factor $\theta$ are
$A_X(\theta)\otimes1$ and $1\otimes M(\theta)$.
Multiplication of function coefficients gives the pairing and unit.
For a basis $e_j$ of $M_m$, the polynomial law with coordinates
\[
 \gamma_p\!\left(\sum_j f_j e_j\right)
 =\sum_{\sum\alpha_j=p}
 \left(\prod_j f_j^{\alpha_j}\right)
 \prod_j\gamma_{\alpha_j}(e_j)
\]
gives $\rho_p$ by universal linearization.  Evaluation verifies basis
independence, every simplicial compatibility, and the divided-power
identities.  All constructions are inverse because their values at
every simplex are the same pointwise maps.
\end{proof}

\paragraph{A-S09-05 (Known; opened comparisons and their exact input).}
Horel, \emph{Binomial rings and homotopy theory}, arXiv:2211.02349v2,
Theorem 3.7, supplies the model structure on cosimplicial binomial rings;
Theorem 5.4 identifies the derived unit
$Y\to\mathbb R\operatorname{Map}_{\mathsf{cBRing}}(A_Y,\mathbb Z)$
as a weak equivalence for connected nilpotent finite-type $Y$.
His source was downloaded and opened; the respective theorem statements
are at source lines 399--407 and 643--647.  The derived adjunction of
his Section 4, applied to this unit, gives
\[
 \mathbb R\operatorname{Map}_{\mathsf{cBRing}}(A_Y,A_X)
 \simeq\mathbb R\operatorname{Map}_{\mathsf{Spaces}}(X,Y).
 \tag{S09.4}
\]
Here $A_X$ is precisely the ring reconstructed in A-S09-03, whose
binomial operations are the unique pointwise integer-valued ones.
Mandell, \emph{Cochains and homotopy type}, published pages 213--216,
Main Theorem and Theorems 0.1--0.2, supplies the integral
$E_\infty$ comparison.  These opened statements have respectively
$A_X$ and the $E_\infty$ cochains as their input.  Formula (S09.4)
does not contain a strict, weight-preserving map of the de Rham
theories from cote 136.

\paragraph{A-S09-04 (Proved; a raw row-test invariant fails).}
There are admissible integral theories $F$ and $A$, with flat terms
and support $q\leq w$, and actual coefficient-preserving PD DGA
morphisms
\[
 F\xrightarrow{\iota}A\xrightarrow{\pi}F,
 \qquad\pi\iota=1_F,
 \tag{S09.5}
\]
whose raw additive row-test functors differ already at a disk.
Consequently the assignment
\[
 A\longmapsto\left[
 L\longmapsto\operatorname{Hom}_{s\mathsf{Ab}}(A^{w,q},L)
 \right]
 \tag{S09.6}
\]
does not send every morphism of admissible theories to an isomorphism
of test functors.  Its direct application cannot give a separating
invariant in the localization inverting all those morphisms.

\begin{proof}
Let $\Phi_n=\mathbb Z^{n+1}$ with marked constant vector $T$, and
$\Psi=\Phi/\mathbb ZT$.  Put
\[
 F^{w,q}=\Gamma^{w-q}(\Phi)\otimes\Lambda^q(\Psi),
 \qquad w\geq q,
\]
with the usual PD de Rham differential, and set the other terms to
zero.  Splitting $\Phi_n=\mathbb ZT\oplus\mathbb Z^n$ identifies
the simplexwise complex with $\Gamma(T)$ tensored with $n$ elementary
PD de Rham complexes.  In one variable the weight-$r>0$ differential
$x^{[r]}\mapsto x^{[r-1]}dx$ is an identity between rank-one groups.
Tensoring these split complexes proves the vertical coefficient
resolution in every weight.
The simplicial group $\Phi=R_0$ is split acyclic by A-S09-02.
Its positive divided powers are simplicially contractible: applying
a reduced functor to the maps of a simplicial contraction preserves
all the defining composition identities.  Tensoring that contraction
with the constant identity homotopy of $\Lambda^q\Psi$ proves
horizontal acyclicity for $w>q$.  Thus $F$ is admissible with
horizontal bound $d+1$ for cochain degrees through $d$.

Let $M^{2r,2r}$ be the constant simplicial group $\mathbb Z e_r$
for each $r\geq1$, and set its other bidegrees and its differential
to zero.  It is an $F$-module through the augmentation
$\epsilon:F\to\mathbb Z$.  Define $A=F\oplus M$ by
\[
 (a,m)(b,n)=(ab,\epsilon(a)n+\epsilon(b)m),\qquad
 d(a,m)=(da,0).
 \tag{S09.7}
\]
In particular $M^2=0$ and the positive-weight ideal of $F$ kills $M$.
For even elements in the augmentation ideal define
$\gamma_j(a,m)=(\gamma_j(a),0)$ for $j\geq2$, and use
$\gamma_1(a,m)=(a,m)$ and $\gamma_0=1$.
These operations satisfy the PD sum, product, scalar and iteration
identities: all mixed terms involving $M$ vanish, while the remaining
terms are exactly the identities in $F$.  For $j=2$ the differential
rule uses $(a,m)(da,0)=(a\,da,0)$; for $j>2$ it is directly the
rule in $F$.  All added generators have even cochain degree, so the
alternating condition is inherited.  These checks apply integrally
and remain valid after every base change, including characteristics
two and three.  The inclusions and projections in (S09.5) preserve
these formulas, both gradings, the coefficient map and the augmentation.

For a fixed $d$, all added terms at weights $w>d$ have cochain degree
$q=w>d$.  The horizontal rows and vertical complex through degree
$d$ therefore agree with $F$ at those weights, except that verification
of the differential out of degree $d$ may use degree $d+1$; taking
$w>d+1$ handles that endpoint as well.  Thus the eventual conditions
hold uniformly in the simplicial degree.  All terms are free and
the support remains $q\leq w$.

At $(w,q)=(2,2)$, the simplex-zero value of $F$ is zero because
$\Psi_0=0$, whereas that of $A$ is $\mathbb Z$.
For every simplicial group $V$, the disk adjunction gives
\[
 \operatorname{Hom}_{s\mathsf{Ab}}(V,D_0)
 =\operatorname{Hom}_{\mathbb Z}((NV)_0,\mathbb Z).
\]
Hence the functor (S09.6) has value zero on $D_0$ for $F$ and value
$\mathbb Z$ for $A$.  These test functors are relatively left exact,
so A-S09-01 applies to them, but they are not isomorphic.  On the
other hand $\pi$ is inverted in the requested localization.  Any
functor out of that localization must send $\pi$ to an isomorphism;
(S09.6) fails this necessary condition, proving the assertion.
\end{proof}

\paragraph{Unsettled step retained explicitly.}
A-S09-01--03 verify the supplied tools.  A-S09-04 checks one exact
attempt to turn their raw test data into an invariant.  No comparison
has yet been constructed from an arbitrary admissible strict PD
theory to the image of $X\mapsto F_X$ that retains all weights,
full natural divided-power actions and actual compatible target maps.
In particular these statements do not construct the local completions
of the finite form truncations required by SA08 C34.  The common-source
and terminal-localization questions over $\mathbb Z$ remain open in
this part of the work.
